\documentclass[11pt,a4paper]{article}

\usepackage[utf8]{inputenc}
\usepackage[T1]{fontenc}
\usepackage[margin=1in]{geometry}
\usepackage{amsmath, amssymb, mathtools, bm}
\usepackage{booktabs}
\usepackage{tabularx}
\usepackage{array}
\usepackage{enumitem}
\usepackage{xcolor}
\usepackage{hyperref}
\usepackage{listings}
\usepackage{longtable}

\hypersetup{
  colorlinks=true,
  linkcolor=blue!50!black,
  citecolor=blue!50!black,
  urlcolor=blue!60!black
}

\lstset{
  basicstyle=\ttfamily\footnotesize,
  breaklines=true,
  columns=fullflexible,
  frame=single,
  showstringspaces=false
}

\newcommand{\Utarget}{U_{\mathrm{target}}}
\newcommand{\Ucirc}{U_{\mathrm{circ}}}
\newcommand{\CNOT}{\mathrm{CNOT}}
\newcommand{\STOP}{\langle \mathrm{STOP}\rangle}
\newcommand{\BOS}{\langle \mathrm{BOS}\rangle}
\newcommand{\Tr}{\mathrm{Tr}}
\newcommand{\E}{\mathbb{E}}
\newcommand{\R}{\mathbb{R}}
\newcommand{\C}{\mathbb{C}}
\newcommand{\softmax}{\mathrm{softmax}}

\title{\textbf{Hybrid-Action Transformer RL for Unitary Synthesis:\\
Literature Review, Code Review, Benchmark Diagnosis,\\
and a Roadmap Toward Publishable State-of-the-Art Results}}
\author{Prepared for the \texttt{gqe-torch} project}
\date{30 April 2026}

\begin{document}
\maketitle

\begin{abstract}
This report reviews the current \texttt{gqe-torch} project for the task of
unitary synthesis with a hybrid-action transformer trained by a reinforcement
learning reward. The code already contains several strong ideas: a GPT-style
autoregressive policy, hybrid discrete/continuous actions, a Pareto archive,
inner and post-training angle refinement, structural simplification, and a
QASER-inspired reward. However, the current formulation is not yet publishable
as a state-of-the-art unitary synthesizer. The most important issue is the
fixed heuristic \texttt{max\_gates\_count}: it is used as both an engineering
limit and a scientific prior, which prevents the policy from learning the
circuit length as part of the task. There are also major benchmark fairness
issues, missing baselines, target-conditioning limitations, and several
training-objective problems.

The main recommendation is to reformulate the method as a \emph{target-
conditioned, variable-horizon, multi-objective synthesis agent} whose actions
operate at the level of expressive skeletons and local synthesis blocks rather
than single primitive gates only. A publishable path is to combine: (i)
variable-length STOP/hazard modeling with an engineering-only cap; (ii)
target/residual conditioning through an encoder or residual-unitary features;
(iii) skeleton-first search inspired by recent gradient-descent results for
generic unitaries; (iv) Rotosolve/closed-form or linear-trace refinement for
angles; (v) Pareto-conditioned or GFlowNet-style training for diverse compact
solutions; and (vi) a benchmark suite including exact Qiskit, Qiskit AQC,
BQSKit QSearch/LEAP/QFAST, and matched-fidelity approximate baselines.
\end{abstract}

\tableofcontents
\newpage

% ---------------------------------------------------------------------------
\section{Problem Definition}
% ---------------------------------------------------------------------------

\subsection{Unitary synthesis}

For $n$ qubits, let $d=2^n$ and $\Utarget \in U(d)$ be the target unitary.
Given a native gate alphabet $\mathcal{G}$, the unitary synthesis problem is
to find a sequence of gates and continuous parameters
\[
  (g_1,\theta_1),\ldots,(g_T,\theta_T), \qquad g_t \in \mathcal{G},
\]
such that
\[
  \Ucirc(\bm{g},\bm{\theta})
  =
  G_T(\theta_T)\cdots G_2(\theta_2)G_1(\theta_1)
\]
approximates $\Utarget$ while using as few expensive operations as possible.
The project uses process fidelity
\[
  F(\Utarget,\Ucirc)
  =
  \frac{|\Tr(\Utarget^\dagger \Ucirc)|^2}{d^2},
\]
so the natural fidelity loss is $1-F$. The current refinement code also uses
the linear-trace loss
\[
  L_{\mathrm{lin}}(\Utarget,\Ucirc)
  =
  1-\frac{|\Tr(\Utarget^\dagger \Ucirc)|}{d},
\]
which has stronger gradients near $F \approx 1$ than $1-F$.

\subsection{Objectives}

For NISQ and early fault-tolerant circuits, the relevant objectives are not
only fidelity. A useful synthesis result should report at least:
\begin{itemize}[leftmargin=2em]
  \item process fidelity $F$ or infidelity $1-F$;
  \item two-qubit gate count, especially CX/CNOT count;
  \item circuit depth under the target hardware connectivity and gate durations;
  \item total primitive gate count;
  \item classical compilation/training time;
  \item success probability across target instances and random seeds.
\end{itemize}
This is a constrained multi-objective problem. The clean comparison is usually:
\[
  \min (C,D,G)
  \quad\text{subject to}\quad
  F \ge 1-\epsilon,
\]
or a Pareto-front comparison at fixed fidelity thresholds. Reporting a single
highest-fidelity circuit and then comparing its structure against an exact
compiler is not sufficient.

\subsection{Lower-bound scale}

Shende, Markov, and Bullock proved a worst-case lower bound for arbitrary
$n$-qubit unitaries over CNOT plus arbitrary one-qubit gates:
\[
  C_{\min}^{\mathrm{worst}}(n)
  \ge
  \left\lceil \frac{4^n-3n-1}{4} \right\rceil .
\]
This is a useful sanity check for the project because it shows that monolithic
arbitrary dense-unitary synthesis is exponentially long.

\begin{center}
\begin{tabular}{rrrr}
\toprule
$n$ & $d=2^n$ & CNOT lower bound & rough $5 \times$ CNOT gate budget \\
\midrule
2 & 4   & 3     & 15 \\
3 & 8   & 14    & 70 \\
4 & 16  & 61    & 305 \\
5 & 32  & 252   & 1260 \\
6 & 64  & 1020  & 5100 \\
7 & 128 & 4091  & 20455 \\
8 & 256 & 16378 & 81890 \\
\bottomrule
\end{tabular}
\end{center}

The current transformer has a hard positional limit of 1024 tokens and the
configuration function caps \texttt{max\_gates\_count} at 1023
(\texttt{config.py:41--76}). Therefore the current monolithic architecture
cannot represent generic 6--8 qubit exact circuits, let alone learn them.
Beating Qiskit up to 8 qubits is feasible only if the benchmark targets are
structured, if the method is hierarchical/blockwise, or if the task is
approximate at a specified fidelity tolerance.

% ---------------------------------------------------------------------------
\section{Current Project Setup}
% ---------------------------------------------------------------------------

\subsection{Model and action space}

The policy is a decoder-only transformer with a discrete token head and a
continuous scalar-angle head (\texttt{policy.py}). The vocabulary is
\[
  \{\BOS,\STOP\} \cup \{\text{rotation tokens, SX tokens, CNOT tokens}\}.
\]
For each decision step, the policy emits:
\[
  a_t=(a_t^d,a_t^c),\qquad a_t^d \in \mathcal{V},\quad a_t^c\in\R.
\]
The continuous action is used only when the discrete token is a parametric
rotation. The discrete sample is drawn from
\[
  \pi(a_t^d=v\mid s_t)
  =
  \softmax(-\beta \ell_t)_v,
\]
where the code treats logits as energies: lower logits are more likely
(\texttt{policy.py:375--377}). The continuous angle is sampled from a
Gaussian with predicted mean and clipped log-standard-deviation
(\texttt{policy.py:337--344}, \texttt{policy.py:389--393}).

\subsection{Circuit evaluator}

\texttt{CircuitEvaluator} constructs the circuit unitary in JAX by scanning
over token and angle arrays of fixed length. STOP, BOS, and padding tokens
are no-ops (\texttt{circuit.py:309--388}). The evaluator uses process fidelity
for rewards and linear-trace cost for refinement.

\subsection{Training loop}

Each epoch performs:
\begin{enumerate}[leftmargin=2em]
  \item sample \texttt{num\_samples} fixed-length token/angle sequences
        with a KV-cached transformer rollout;
  \item evaluate raw fidelities;
  \item optionally refine each rollout's angles with Adam
        (\texttt{trainer.py:492--507});
  \item compute the QASER-style scalar reward
        (\texttt{trainer.py:591--635});
  \item store samples in a replay buffer;
  \item run PPO/GRPO-style updates using stored log probabilities;
  \item update a Pareto archive of non-dominated circuits.
\end{enumerate}
Post-training, archive entries are refined again with longer Adam runs and
optional closed-form sweep polish.

\subsection{Benchmark setup observed in \texttt{results/benchmarks}}

The benchmark script trains a fresh policy per target, compiles the same dense
target with Qiskit's default unitary synthesis at \texttt{optimization\_level=3},
and writes JSONL rows. Current complete benchmark files contain 50 targets for
2q, 3q, and 4q brickwork-Haar target distributions. The latest complete 3q
run is strong on depth and total gate count. The 4q run is high fidelity but
not exact and loses badly on CNOT count.

\begin{center}
\small
\begin{tabular}{llrrrrrr}
\toprule
file & $n$ & mean $F$ & success $F\ge0.99$ & mean CX GQE & mean CX Qiskit & mean depth GQE & mean depth Qiskit \\
\midrule
\texttt{20260428\_122555} & 2 & 0.9998 & 50/50 & 2.96 & 3  & 14.18  & 15 \\
\texttt{20260428\_141635} & 2 & 0.9991 & 49/50 & 2.74 & 3  & 14.94  & 15 \\
\texttt{20260428\_144945} & 3 & 0.9995 & 50/50 & 18.74 & 19 & 51.96  & 62 \\
\texttt{20260429\_135647} & 3 & 0.9999 & 50/50 & 19.08 & 19 & 50.84  & 62 \\
\texttt{20260429\_021633} & 4 & 0.9963 & 50/50 & 119.34 & 95 & 199.68 & 272 \\
\bottomrule
\end{tabular}
\end{center}

The current results show a promising structural trade-off: GQE often produces
shallower circuits than Qiskit. However, for 4q it uses about 24 more CNOTs on
average and reaches only about 0.996 mean fidelity. This is not yet a
state-of-the-art result.

% ---------------------------------------------------------------------------
\section{Literature Review}
% ---------------------------------------------------------------------------

\subsection{Exact and constructive unitary synthesis}

\paragraph{CSD/QSD and generic lower bounds.}
Classical exact synthesis decomposes a dense unitary using algebraic
factorizations such as the cosine-sine decomposition (CSD) and Quantum
Shannon Decomposition (QSD). Shende, Bullock, and Markov's synthesis work
improved the asymptotic CNOT count for arbitrary $n$-qubit unitaries and
established the lower-bound scale used above \cite{shende2006synthesis}.
These algorithms are reliable but exponential in $n$ and become impractical
or very deep beyond small blocks.

\paragraph{Two-qubit optimal synthesis.}
For two-qubit unitaries, KAK/Cartan/Weyl decomposition gives optimal
decompositions in terms of a two-qubit basis gate. Qiskit's synthesis module
exposes \texttt{TwoQubitBasisDecomposer} and a \texttt{two\_qubit\_cnot\_decompose}
entry point for CX-based two-qubit decomposition \cite{qiskit_synthesis}.
This matters because most scalable synthesis systems reduce large circuits
to small blocks and instantiate the bottom blocks with KAK or equivalent
two-qubit routines. The three-CNOT worst-case optimality of arbitrary
two-qubit unitaries is the canonical small-block reference point
\cite{shende2004minimal}.

\subsection{Modern synthesis toolchains}

\paragraph{Qiskit.}
Qiskit supports a unitary synthesis plugin interface. Current documentation
lists built-in unitary synthesis plugins including \texttt{default},
\texttt{aqc}, \texttt{sk}, and others depending on installation
\cite{qiskit_plugins}. The \texttt{UnitarySynthesis} pass can do exact or
approximate two-qubit synthesis depending on \texttt{approximation\_degree}
and target information \cite{qiskit_unitary_synthesis}. The AQC module
implements an approximate compiler using CNOT networks and continuous
optimization, with L-BFGS as the default optimizer in the implementation
\cite{qiskit_aqc,madden2021aqc}. Therefore a fair benchmark against
``Qiskit'' should include at least default exact synthesis, AQC with a
depth/CNOT sweep, fixed seeds, and matched fidelity thresholds.

\paragraph{BQSKit.}
BQSKit is a strong modern baseline because it combines partitioning,
synthesis, and instantiation. Its documentation describes QSearch as
optimal-depth and topology-aware up to about four qubits, LEAP as a higher
quality extension up to about six qubits, and QFAST as a scalable approximate
approach up to about eight qubits \cite{bqskit_home,bqskit_synthesis}. This
is directly relevant to the stated goal of beating Qiskit up to 8-qubit
unitaries: BQSKit is likely a stronger baseline than Qiskit's default dense
\texttt{UnitaryGate} synthesis for structured targets. QFAST is especially
relevant because it uses a hierarchical continuous circuit space rather than
an exhaustive primitive-gate tree \cite{qfast2020}.

\paragraph{Recent gradient-descent result for generic unitaries.}
Gomathi and Meiburg (2026) report that simple gradient descent reliably finds
depth- and gate-optimal circuits for generic unitaries when the CNOT skeleton
is chosen to avoid parameter-deficient layouts \cite{gomathi2026gd}. This is
highly relevant: it suggests the hard part is not always continuous-angle
optimization, but choosing an expressive skeleton. Your current code already
contains a pair-repetition filter inspired by this line of work, but the
policy still samples primitive gate sequences rather than structured
skeleton families.

\subsection{RL and generative models for quantum circuit generation}

\paragraph{Deep RL quantum compiling.}
Moro et al. formulate quantum compiling as an RL environment that builds a
circuit one gate at a time and uses fidelity-based rewards; they demonstrate
learned single-qubit compiling policies \cite{moro2021drl}. Foelsel et al.
use deep RL for architecture-aware quantum circuit optimization and report
average depth and gate-count reductions on 12-qubit random circuits
\cite{fosel2021rl}. These works support the RL framing, but they also show
that benchmark design and hardware-aware observations are central.

\paragraph{Hybrid discrete/continuous action RL.}
The project's hybrid token-plus-angle action belongs to the parameterized
action-space RL family. Hybrid Actor-Critic and H-PPO decompose structured
hybrid actions into coordinated discrete and continuous policies
\cite{fan2019hppo}. Compared with those methods, the current code uses one
sequence-level PPO ratio and a masked continuous log-prob. This is workable,
but the continuous head and entropy should be conditioned more carefully on
the chosen discrete action.

\paragraph{GQE and transformers.}
The Generative Quantum Eigensolver (GQE) reframes circuit construction as
sequence generation with a classical generative model. Nakaji et al. introduce
GQE and a GPT-style implementation for ground-state search
\cite{nakaji2024gqe,nvidia2024gqe}. Conditional-GQE extends this idea with an encoder-
decoder transformer so that the generator is conditioned on the problem
instance, reaching strong generalization on combinatorial optimization
instances up to 10 qubits \cite{minami2025conditional}. This is one of the
most important lessons for this project: training a new unconditioned
transformer from scratch for every target is not an attractive compiler story.

\paragraph{Preference and diversity objectives.}
GQE variants increasingly use preference-based training. Persistent-DPO was
proposed to address limitations of direct preference optimization in GQE
\cite{nakamura2025pdpo}. GFlowNets are also relevant because they learn to
sample diverse high-reward compositional objects rather than collapsing to one
mode \cite{zhang2022gflownet}. Since unitary synthesis has many
non-dominated circuits with different CNOT/depth trade-offs, a GFlowNet-style
or preference-conditioned archive objective fits the problem better than a
single scalar reward.

\subsection{Continuous parameter optimization}

The current code uses Adam refinement and optional closed-form sweep polish.
This aligns with the broader lesson that discrete structure search and
continuous instantiation should be separated. Rotosolve-style coordinate
updates and the sweep update already implemented in \texttt{polish.py} are
especially appealing for rotation gates because each single-axis rotation has
a sinusoidal dependence on the objective when all other parameters are fixed
\cite{ostaszewski2021structure}. The project should treat this not merely as
post-processing, but as part of the inner training signal for skeleton
evaluation.

% ---------------------------------------------------------------------------
\section{Critical Flaws and Risks in the Current Code}
% ---------------------------------------------------------------------------

\subsection{The fixed \texttt{max\_gates\_count} is the central design flaw}

The project currently sets a hard maximum sequence length through
\texttt{default\_max\_gates\_count} (\texttt{config.py:41--76}). Benchmarking
also recomputes this heuristic for every qubit count
(\texttt{benchmark.py:83--107}). This creates three problems.

\paragraph{Scientific bias.}
The policy is supposed to learn the correct number of gates from reward.
Instead, the code pre-selects a horizon based on a rough estimate of generic
Qiskit-like decompositions. This can make an experiment look good or bad for
the wrong reason. If the cap is too small, high-fidelity solutions are
unrepresentable. If it is too large, rollouts waste compute and the reward
must fight a large space of redundant STOP-padded sequences.

\paragraph{No true variable-horizon computation.}
The policy has a STOP token, but generation still scans the entire fixed
sequence length even after all samples have stopped (\texttt{policy.py:347--422}).
Thus \texttt{max\_gates\_count} remains a runtime and memory bottleneck.

\paragraph{Incompatibility with generic 8q synthesis.}
The 1023-token cap is below even rough generic budgets for $n\ge 6$. For
8-qubit dense arbitrary unitaries the CNOT lower bound is 16378, so an
8-qubit state-of-the-art claim must be restricted to structured targets,
approximate synthesis, or a hierarchical blockwise architecture.

\subsection{The policy is not target-conditioned}

\texttt{HybridPolicy} receives previous tokens only. The target unitary is not
an input to the transformer; instead, a fresh policy is trained from scratch
for each target. This is the biggest scalability and publishability issue
after the fixed horizon. A compiler should amortize knowledge across targets
or at least condition on the current target/residual. Current training is more
like black-box architecture search for one matrix at a time.

Recommended target-conditioning options:
\begin{itemize}[leftmargin=2em]
  \item feed an embedding of $\Utarget$ to the transformer through cross-attention;
  \item feed the residual $R_t = \Utarget U_t^\dagger$ or features of it at each step;
  \item encode local block residuals in a hierarchical synthesizer;
  \item condition on scalarization weights, fidelity tolerance, and hardware constraints.
\end{itemize}

\subsection{Primitive-token search is too low-level}

The current vocabulary samples single rotations, SX gates, and ordered CNOTs.
This makes the transformer responsible for discovering elementary compiler
grammar from scratch. Literature and existing tools suggest a better level of
abstraction:
\begin{itemize}[leftmargin=2em]
  \item sample CNOT skeletons and instantiate all one-qubit layers with a
        continuous optimizer;
  \item use KAK/SU(4) macro blocks on 2-qubit neighborhoods;
  \item use QFAST-style $m$-qubit blocks and refine them hierarchically;
  \item use hardware-native entangling layers or one-factorization patterns
        as skeleton tokens.
\end{itemize}
This would reduce the search horizon and align the method with both BQSKit
and the 2026 gradient-descent result on generic-unitary skeletons.

\subsection{Reward design is unstable and not lexicographic enough}

The QASER-style reward in \texttt{trainer.py:611--635} is inspired by
QASER's attempt to balance depth, two-qubit count, and accuracy in a
reinforcement-learning reward \cite{qaser2025}. It uses running maxima
$M_D,M_C,M_G$ from observed batches:
\[
  \mathrm{base}
  =
  w_D\frac{M_D}{D+1}
  +
  w_C\frac{M_C}{C+1}
  +
  w_G\frac{M_G}{G+1},
  \qquad
  R=\mathrm{base}^{F}\cdot \text{log multiplier}-1.
\]
The maxima are nonstationary and path-dependent. A single early batch with
very long circuits changes the reward scale for the rest of training. Also,
the reward is a scalarization of objectives that are better handled
lexicographically or by constraints.

For unitary synthesis, a more robust objective is:
\[
  R =
  -\alpha \log(1-F+\epsilon)
  - \lambda_C C - \lambda_D D - \lambda_G G
  - \lambda_{\mathrm{nostop}}\mathbf{1}\{\text{no STOP}\},
\]
with $\lambda$ terms activated or increased only after a target fidelity band
has been reached. Even better, sample scalarization weights during training
and train a Pareto-conditioned policy, so the model learns the whole trade-off
front instead of one hand-tuned reward.

\subsection{GRPO/PPO is used in an off-policy replay setting}

Advantages are computed inside each PPO minibatch using
\texttt{grpo\_advantages(costs)} (\texttt{trainer.py:421}), but the minibatch
is drawn from a replay buffer, not necessarily the original rollout group. This
is not the same as group-relative normalization over samples generated by the
same behavior policy. It can change signs and scales of advantages depending
on arbitrary buffer composition. A cleaner implementation would store the
rollout-group mean/std at collection time, store normalized advantages, and
either avoid replay or use explicit off-policy corrections.

\subsection{Temperature schedule is applied at the wrong time}

\texttt{collect\_rollout} samples with the current beta and then updates the
scheduler before PPO training (\texttt{trainer.py:479--544}). The subsequent
PPO step uses the updated beta (\texttt{trainer.py:694--695}), while the old
log probabilities were produced under the previous beta. Importance ratios do
account for old log probabilities, but beta is an exogenous policy parameter.
Changing it between data collection and policy update makes the update optimize
a different sampling distribution from the one that produced the data. The
simpler rule is: use the same beta for rollout and PPO updates, then update the
scheduler after training on that rollout.

\subsection{Continuous angle distribution ignores periodicity}

Angles are $2\pi$-periodic, but the policy density is an unbounded Gaussian.
The circuit evaluator treats $\theta$ and $\theta+2\pi k$ as equivalent, while
PPO treats them as different actions with different log probabilities. This
creates avoidable variance and can punish equivalent good actions. A wrapped
normal, von Mises distribution, or explicit modulo canonicalization with the
correct likelihood would better match the physics.

\subsection{Continuous entropy is applied to non-parametric tokens}

The PPO loss masks continuous log-probability contributions for non-parametric
tokens, but the continuous entropy term uses \texttt{valid} only
(\texttt{trainer.py:424--425}), not \texttt{param\_mask \& valid}. Thus the
continuous head is regularized on CNOT/SX/STOP positions where the sampled
angle has no physical meaning. This is not catastrophic, but it wastes
capacity and can inject irrelevant gradients.

\subsection{Pareto archive ignores total gates in dominance}

\texttt{ParetoArchive} tracks fidelity, depth, total gates, and CNOT count, but
dominance only uses fidelity, depth, and CNOT count
(\texttt{pareto.py:100--110}). Total gates are later queried, but they are not
part of non-dominance. This can preserve circuits that are bloated in total
single-qubit count as long as depth and CNOT count are competitive. Since the
reward includes total gates, the archive should either include total gates in
dominance or maintain separate constrained fronts.

\subsection{Best circuit bookkeeping is misleading}

The code comments say \texttt{best\_raw\_fidelity} tracks the highest-fidelity
raw rollout, but when inner refinement is enabled, \texttt{fids} are already
post-inner-refinement (\texttt{trainer.py:497--511}, \texttt{trainer.py:688--692}).
This makes logs and benchmark interpretation less clear. Also,
\texttt{early\_stop} checks the fidelity associated with the best-cost circuit
(\texttt{trainer.py:681--707}), not the best-fidelity circuit. A high-fidelity
sample could appear without triggering early stop.

\subsection{Config and metadata are incomplete}

Benchmark metadata writes the base config, not the per-target overridden config.
For example, some 2q benchmark metadata report \texttt{max\_gates\_count=458}
even though the benchmark rows are 2q targets. This happens because metadata is
written from \texttt{base\_cfg}, while each target uses an overridden config.
The metadata also omits several important fields such as
\texttt{brickwork\_depth} in the config snapshot, some refinement knobs, and
the git commit/dependency versions.

% ---------------------------------------------------------------------------
\section{Benchmarking Problems}
% ---------------------------------------------------------------------------

\subsection{The current Qiskit comparison is too narrow}

\texttt{benchmark.py} uses:
\begin{lstlisting}[language=Python]
compiled = transpile(qc, basis_gates=basis, optimization_level=3)
\end{lstlisting}
and assumes Qiskit fidelity is exactly 1
(\texttt{benchmark.py:114--128}). This is a valid first baseline but not a
publishable benchmark. A fair suite should include:
\begin{itemize}[leftmargin=2em]
  \item Qiskit default exact unitary synthesis;
  \item Qiskit AQC with depth/CNOT sweeps and fixed seeds;
  \item Qiskit approximate synthesis at matched target fidelity;
  \item BQSKit QSearch/LEAP/QFAST where applicable;
  \item possibly TKET and/or pyZX-style optimization for circuit inputs;
  \item multiple \texttt{seed\_transpiler} values for stochastic passes.
\end{itemize}

\subsection{The selected GQE circuit is not the right constrained optimum}

\texttt{\_gqe\_best} selects \texttt{arc.best\_by\_fidelity()}
(\texttt{benchmark.py:131--146}). If the archive contains a lower-CNOT circuit
with $F\ge0.99$, the benchmark ignores it. Conversely, if the best-fidelity
GQE circuit has $F=0.996$ and Qiskit is exact, comparing structure directly is
not apples-to-apples. The benchmark should report:
\begin{itemize}[leftmargin=2em]
  \item best fidelity;
  \item minimum CNOT count at $F\ge0.99$, $0.999$, $0.9999$, and exact tolerance;
  \item minimum depth at the same thresholds;
  \item success rate at each threshold;
  \item median paired deltas and confidence intervals.
\end{itemize}

\subsection{Target distributions need to be separated}

The current benchmark uses brickwork-Haar targets. These are important, but
they are not the same as arbitrary Haar-random dense unitaries. For a
publishable paper, separate at least four categories:
\begin{enumerate}[leftmargin=2em]
  \item 2--4q Haar-random dense unitaries, where exact baselines are meaningful;
  \item 2--8q brickwork/local random circuits, where locality recovery matters;
  \item structured algorithmic blocks, such as QFT, arithmetic, chemistry, and
        Hamiltonian simulation fragments;
  \item hardware-constrained targets with a fixed coupling map and native basis.
\end{enumerate}
If the claim is ``beating Qiskit up to 8 qubits,'' state whether Qiskit receives
a dense \texttt{UnitaryGate} or the original circuit. Beating Qiskit on dense
matrices generated from hidden shallow circuits is interesting, but it is a
different claim from beating Qiskit when both methods receive the same circuit.

\subsection{Statistical reporting is missing}

The plots show point estimates. A top-conference paper needs paired statistics:
bootstrap confidence intervals, Wilcoxon signed-rank tests for paired deltas,
success-rate confidence intervals, and ablations. Since GQE is stochastic and
Qiskit can be stochastic, each target should have multiple seeds or a clearly
stated deterministic protocol.

% ---------------------------------------------------------------------------
\section{Fixing the Gate-Count Heuristic}
% ---------------------------------------------------------------------------

\subsection{Separate engineering cap from learned circuit length}

The project should keep a cap $T_{\max}$ only as an engineering safety limit,
not as the scientific gate budget. The learned object should be variable length:
\[
  p_\phi(\tau, a_{1:\tau}\mid \Utarget, \epsilon, \bm{w}),
  \qquad
  \tau = \min\{t:a_t=\STOP\}.
\]
The reward should use the actual non-noop gate count $G$, CNOT count $C$, and
depth $D$, not the cap. No-STOP samples should receive an explicit terminal
penalty because they failed to decide a circuit length.

\subsection{Use a STOP hazard head}

Instead of treating STOP as just another token inside a huge vocabulary, use a
separate Bernoulli hazard head:
\[
  h_t = p_\phi(\mathrm{stop}\mid s_t).
\]
If not stopped, sample a gate token from the non-STOP vocabulary. This makes
length modeling explicit and gives cleaner entropy/KL control over termination.
The policy can still be implemented with a STOP token internally, but the loss
should expose the hazard term separately.

\subsection{Use iterative deepening or adaptive caps}

For each target, start with a small cap derived from a lower-bound estimate or
from the best known baseline. If no circuit reaches the fidelity threshold,
increase the cap. This mirrors exact-search practice and avoids wasting 1023
steps on easy 2q/3q targets. In pseudocode:
\[
  T_{\max}^{(0)} < T_{\max}^{(1)} < \cdots
\]
and train/evaluate until a front with $F\ge1-\epsilon$ is found. The reported
length remains learned; the cap is only a search envelope.

\subsection{Train with budget conditioning}

Sample a budget token $b$ or scalarization weights $\bm{w}$ at the start of an
episode:
\[
  p_\phi(a_{1:\tau}\mid \Utarget, b,\bm{w}).
\]
Sometimes ask for high fidelity regardless of length; sometimes ask for CNOT
minimality at fixed fidelity. This teaches the transformer the Pareto front and
prevents a single reward from baking in one trade-off.

% ---------------------------------------------------------------------------
\section{Proposed Publishable Research Direction}
% ---------------------------------------------------------------------------

\subsection{Core thesis}

A strong publishable contribution would be:

\begin{quote}
\emph{A target-conditioned generative synthesis agent that learns a distribution
over variable-length circuit skeletons and uses differentiable/closed-form
angle instantiation to produce a Pareto front of high-fidelity, low-CNOT,
low-depth circuits.}
\end{quote}

This is more novel and defensible than ``PPO over primitive gate tokens with a
hand-tuned max length.'' It connects GQE/conditional transformers, hybrid
actions, generic-unitary skeleton results, and multi-objective synthesis.

\subsection{Architecture}

\paragraph{Encoder.}
Encode the target or residual. For $n\le4$, the dense unitary can be embedded
directly through real/imaginary channels, singular/invariant features, or a
small transformer over matrix rows/columns. For $n>4$, use local block
residuals, Pauli-transfer features, tensor-network features, or features from
a partitioning pass.

\paragraph{Decoder.}
Use an autoregressive transformer decoder with cross-attention to the target
embedding. Add conditioning tokens for hardware graph, fidelity threshold,
and objective weights.

\paragraph{Actions.}
Use hierarchical actions:
\begin{itemize}[leftmargin=2em]
  \item high level: choose a block, skeleton motif, or entangling edge pattern;
  \item mid level: choose CNOT/SU(4)/KAK macro placement;
  \item low level: emit rotation-axis tokens only where a skeleton requires
        free parameters.
\end{itemize}
Primitive gate emission can remain as a fallback mode for fine repair.

\paragraph{Angle instantiation.}
For each sampled skeleton, instantiate parameters with a fast inner solver:
linear-trace Adam for warm starts, followed by one or two sweep/Rotosolve
passes. The RL reward should evaluate the instantiated skeleton, not raw random
angles only.

\subsection{Training objective}

Use a constrained or Pareto-conditioned objective:
\[
  R_{\bm{w},\epsilon}
  =
  -\log(1-F+\delta)
  - \mathbf{1}\{F\ge 1-\epsilon\}
    (w_C C + w_D D + w_G G)
  - \lambda_{\mathrm{fail}}\mathbf{1}\{\text{no STOP}\}.
\]
Alternatively, train with archive-derived preferences:
\[
  y^+ \succ y^-
  \quad\text{if}\quad
  y^+ \text{ dominates } y^- \text{ or is better under sampled } \bm{w}.
\]
Then apply DPO/Persistent-DPO-style updates or GFlowNet flow matching to
sample diverse high-quality circuits instead of collapsing to one mode.

\subsection{Why this can beat Qiskit}

Qiskit's default exact dense-unitary path is general and reliable but not
specialized to hidden local structure. A target-conditioned generative model
trained on brickwork/local random circuits can learn locality priors and
recover shallow structures that dense decompositions miss. Qiskit AQC and
BQSKit are much stronger baselines, so the proposed method should compete
where it has a real advantage:
\begin{itemize}[leftmargin=2em]
  \item amortized generation across many related targets;
  \item fast production of a whole Pareto front;
  \item learned priors over hardware-native skeletons;
  \item variable-horizon search that can stop early when the target is simple;
  \item gradient/sweep inner solvers that make skeleton evaluation reliable.
\end{itemize}

\subsection{A realistic 8-qubit plan}

Do not try to synthesize arbitrary 8-qubit Haar matrices monolithically with a
1024-token transformer. Instead:
\begin{enumerate}[leftmargin=2em]
  \item define 8q benchmark targets as structured local/brickwork circuits and
        state that clearly;
  \item use a block partitioner to decompose the 8q unitary or circuit into
        2--4q subproblems;
  \item train the model to synthesize 2--4q blocks and stitch them with a
        QFAST/BQSKit-like hierarchy;
  \item compare against Qiskit dense, Qiskit AQC, and BQSKit QFAST at matched
        fidelity;
  \item report both dense-matrix input and original-circuit input settings.
\end{enumerate}

% ---------------------------------------------------------------------------
\section{Concrete Engineering Recommendations}
% ---------------------------------------------------------------------------

\subsection{High-priority changes}

\begin{enumerate}[leftmargin=2em]
  \item Replace \texttt{max\_gates\_count} as a task heuristic with an
        engineering cap plus learned STOP/hazard termination.
  \item Add target conditioning. At minimum, condition on a fixed embedding of
        $\Utarget$ for small $n$.
  \item Move scheduler updates to after PPO training for the rollout that used
        the current beta.
  \item Store rollout-group normalized advantages rather than recomputing GRPO
        advantages on replay minibatches.
  \item Mask continuous entropy with \texttt{param\_mask \& valid}.
  \item Add periodic angle distributions or modulo-aware likelihood.
  \item Include total gate count in Pareto dominance or maintain separate
        fronts for $(F,C,D)$ and $(F,C,D,G)$.
  \item Make benchmark selection constrained: min CNOT/depth at fixed fidelity
        thresholds, not best fidelity only.
  \item Verify and record Qiskit fidelity instead of assuming it.
  \item Add Qiskit AQC and BQSKit baselines.
\end{enumerate}

\subsection{Medium-priority changes}

\begin{enumerate}[leftmargin=2em]
  \item Add a skeleton-only mode where the policy emits CNOT edges/layers and
        the optimizer supplies one-qubit rotations.
  \item Add macro actions for two-qubit SU(4) blocks, later decomposed by KAK.
  \item Implement iterative deepening over caps.
  \item Log git commit, package versions, device, actual per-run config, and
        per-phase timing.
  \item Add bootstrap confidence intervals and paired significance tests to
        benchmark reports.
  \item Add tests for Qiskit fidelity verification, Pareto constrained
        selection, reward monotonicity, and STOP/no-STOP behavior.
\end{enumerate}

\subsection{Ablation matrix for a paper}

\begin{center}
\small
\begin{tabularx}{\textwidth}{lX}
\toprule
Ablation & Question answered \\
\midrule
fixed cap vs learned STOP & Does variable horizon improve compactness and compute? \\
unconditioned vs target-conditioned & Does the model amortize across targets? \\
primitive gates vs skeleton/macro actions & Does abstraction reduce search difficulty? \\
raw RL angles vs inner refinement & Is angle optimization required for reliable reward? \\
Adam vs sweep/Rotosolve vs hybrid & Which instantiator gives best fidelity per compute? \\
scalar reward vs Pareto-conditioned reward & Does the model recover better trade-off fronts? \\
QASER vs constrained log-infidelity reward & Which reward avoids low-fidelity compact traps? \\
Qiskit exact vs Qiskit AQC vs BQSKit & Is the claimed baseline actually state of the art? \\
\bottomrule
\end{tabularx}
\end{center}

% ---------------------------------------------------------------------------
\section{Conclusion}
% ---------------------------------------------------------------------------

The current project is promising but not yet publishable as a state-of-the-art
unitary synthesis method. The benchmark results already show useful behavior:
on 3q brickwork-Haar targets, the method reaches very high fidelity and often
beats Qiskit on depth and total gates. The 4q results show the current limit:
the model finds shallow high-fidelity circuits but loses substantially on CNOT
count and does not reach exact fidelity.

The first fix is exactly the issue identified in the prompt: remove the
scientific dependence on \texttt{max\_gates\_count}. Let the policy learn
when to stop, while using a cap only as a computational envelope. The deeper
fix is to stop training an unconditioned primitive-token policy from scratch
for each target. A publishable system should be target-conditioned, variable
horizon, skeleton-aware, multi-objective, and benchmarked against stronger
baselines at matched fidelity.

% ---------------------------------------------------------------------------
\begin{thebibliography}{99}

\bibitem{shende2006synthesis}
V. V. Shende, S. S. Bullock, and I. L. Markov.
\newblock Synthesis of Quantum Logic Circuits.
\newblock \emph{IEEE Transactions on Computer-Aided Design of Integrated
Circuits and Systems}, 25(6), 2006.
\newblock \url{https://www.nist.gov/publications/synthesis-quantum-logic-circuits}

\bibitem{shende2004minimal}
V. V. Shende, I. L. Markov, and S. S. Bullock.
\newblock Minimal Universal Two-Qubit Controlled-NOT-Based Circuits.
\newblock \emph{Physical Review A}, 69, 2004.
\newblock \url{https://www.nist.gov/publications/minimal-universal-two-qubit-controlled-not-based-circuits}

\bibitem{qiskit_plugins}
IBM Quantum Documentation.
\newblock Install and use transpiler plugins.
\newblock \url{https://qiskit.qotlabs.org/docs/guides/transpiler-plugins}

\bibitem{qiskit_unitary_synthesis}
IBM Quantum Documentation.
\newblock \texttt{UnitarySynthesis} pass.
\newblock \url{https://docs.quantum.ibm.com/api/qiskit/qiskit.transpiler.passes.UnitarySynthesis}

\bibitem{qiskit_synthesis}
IBM Quantum Documentation.
\newblock Circuit synthesis API.
\newblock \url{https://docs.quantum.ibm.com/api/qiskit/synthesis}

\bibitem{qiskit_aqc}
IBM Quantum Documentation.
\newblock Approximate Quantum Compiler module.
\newblock \url{https://qiskit.qotlabs.org/docs/api/qiskit/qiskit.synthesis.unitary.aqc}

\bibitem{madden2021aqc}
L. Madden and A. Simonetto.
\newblock Best Approximate Quantum Compiling Problems.
\newblock arXiv:2106.05649, 2021.
\newblock \url{https://arxiv.org/abs/2106.05649}

\bibitem{bqskit_home}
Berkeley Quantum Synthesis Toolkit.
\newblock BQSKit overview.
\newblock \url{https://bqskit.lbl.gov/}

\bibitem{bqskit_synthesis}
BQSKit Documentation.
\newblock Quantum Synthesis Explained.
\newblock \url{https://bqskit.readthedocs.io/en/latest/intro/synthesis.html}

\bibitem{qfast2020}
A. Younis.
\newblock QFAST: Quantum Synthesis Using a Hierarchical Continuous Circuit Space.
\newblock UC Berkeley Technical Report, 2020.
\newblock \url{https://www2.eecs.berkeley.edu/Pubs/TechRpts/2020/EECS-2020-53.html}

\bibitem{gomathi2026gd}
J. Gomathi and A. Meiburg.
\newblock Gradient descent reliably finds depth- and gate-optimal circuits for generic unitaries.
\newblock arXiv:2601.03123, 2026.
\newblock \url{https://arxiv.org/abs/2601.03123}

\bibitem{qaser2025}
I. Moflic, A. Paler, and A. Kundu.
\newblock QASER: Breaking the Depth vs. Accuracy Trade-Off for Quantum Architecture Search.
\newblock arXiv:2511.16272, 2025.
\newblock \url{https://arxiv.org/abs/2511.16272}

\bibitem{moro2021drl}
L. Moro, M. G. A. Paris, M. Restelli, and E. Prati.
\newblock Quantum compiling by deep reinforcement learning.
\newblock \emph{Communications Physics}, 4, 178, 2021.
\newblock \url{https://www.nature.com/articles/s42005-021-00684-3}

\bibitem{fosel2021rl}
T. Foelsel, M. Y. Niu, F. Marquardt, and L. Li.
\newblock Quantum circuit optimization with deep reinforcement learning.
\newblock arXiv:2103.07585, 2021.
\newblock \url{https://arxiv.org/abs/2103.07585}

\bibitem{fan2019hppo}
Z. Fan, R. Su, W. Zhang, and Y. Yu.
\newblock Hybrid Actor-Critic Reinforcement Learning in Parameterized Action Space.
\newblock IJCAI 2019.
\newblock \url{https://www.ijcai.org/Proceedings/2019/316}

\bibitem{nakaji2024gqe}
K. Nakaji et al.
\newblock The generative quantum eigensolver (GQE) and its application for ground state search.
\newblock arXiv:2401.09253, 2024; revised 2025.
\newblock \url{https://arxiv.org/abs/2401.09253}

\bibitem{minami2025conditional}
S. Minami, K. Nakaji, Y. Suzuki, A. Aspuru-Guzik, and T. Kadowaki.
\newblock Generative quantum combinatorial optimization by means of a novel conditional generative quantum eigensolver.
\newblock \emph{Digital Discovery}, 4, 2229--2243, 2025.
\newblock \url{https://pubs.rsc.org/en/content/articlehtml/2025/dd/d5dd00138b}

\bibitem{nakamura2025pdpo}
J. Nakamura.
\newblock Persistent-DPO: A novel loss function and hybrid learning for generative quantum eigensolver.
\newblock arXiv:2509.08351, 2025.
\newblock \url{https://arxiv.org/abs/2509.08351}

\bibitem{zhang2022gflownet}
D. Zhang, N. Malkin, Z. Liu, A. Volokhova, A. Courville, and Y. Bengio.
\newblock Generative Flow Networks for Discrete Probabilistic Modeling.
\newblock ICML 2022.
\newblock \url{https://proceedings.mlr.press/v162/zhang22v.html}

\bibitem{ostaszewski2021structure}
M. Ostaszewski, E. Grant, and M. Benedetti.
\newblock Structure optimization for parameterized quantum circuits.
\newblock \emph{Quantum}, 5:391, 2021.
\newblock \url{https://quantum-journal.org/papers/q-2021-01-28-391/}

\bibitem{nvidia2024gqe}
M. Wolf, P. Rao, K. Nakaji, and C. Gorgulla.
\newblock Advancing Quantum Algorithm Design with GPTs.
\newblock NVIDIA Technical Blog, 2024.
\newblock \url{https://developer.nvidia.com/blog/advancing-quantum-algorithm-design-with-gpt/}

\end{thebibliography}

\end{document}
